\chapter{Metodología}\label{chap:metodologia}

En este capítulo se debe describir y fundamentar la metodología empleada en la tesis, la cual debe responder directamente a los objetivos e hipótesis planteados en la introducción. Esta descripción debe ser detallada; no se debe presentar un resumen de una sola página. Sin embargo, la explicación debe ser conceptual, sin llegar a los detalles de implementación. Por ejemplo, cualquier referencia al código fuente debe ubicarse en la sección de anexos (Apéndice \ref{chap:AnexoB} en el caso de esta plantilla).

Para apoyar la explicación, suele ser útil incluir un diagrama de flujo, ya sea de la metodología completa o de etapas específicas de esta, como se muestra en la figura \ref{img:Diagrama}. Adicionalmente, en esta sección probablemente se requiera explicar algoritmos, como se muestra en el algoritmo \ref{alg:Algoritmo}.

\begin{figure}[h]
    \centering
    \shorthandoff{>}
    \scalebox{0.7}{\begin{tikzpicture}[
    >=Stealth,
    node distance=0.8cm,
    block/.style={
        rectangle,
        draw=black,
        thick,
        minimum width=2.4cm,
        minimum height=0.9cm,
        align=center,
        font=\small
    },
    decision/.style={
        diamond,
        draw=black,
        thick,
        minimum width=2.4cm,
        minimum height=0.9cm,
        align=center,
        font=\small,
        aspect=2
    },
    terminal/.style={
        rectangle,
        draw=black,
        thick,
        rounded corners=0.4cm,
        minimum width=2.4cm,
        minimum height=0.9cm,
        align=center,
        font=\small
    }
]
 
\node[terminal]  (inicio)    {Inicio};
\node[block,     right=of inicio]    (paso1)    {Paso 1};
\node[block,     right=of paso1]     (paso2)    {Paso 2};
\node[decision,  right=of paso2]     (decision) {Condicion};
\node[block,     right=of decision]  (paso3)    {Paso 3};
\node[terminal,  right=of paso3]     (fin)      {Fin};
 
\draw[->, thick] (inicio)   -- (paso1);
\draw[->, thick] (paso1)    -- (paso2);
\draw[->, thick] (paso2)    -- (decision);
\draw[->, thick] (decision) -- node[above, font=\small] {Si} (paso3);
\draw[->, thick] (paso3)    -- (fin);
 
% Loop back on No
\draw[->, thick] (decision.south)
    -- ++(0,-1.0)
    -| node[below, font=\small, pos=0.25] {No} (paso1.south);
 
\end{tikzpicture}}
    \shorthandon{>}
    \caption{Diagrama de flujo de la metodología.}
    \label{img:Diagrama}
\end{figure}

\begin{algorithm}[h]
\caption{Nombre del algoritmo.}
\label{alg:Algoritmo}
\KwIn{Datos de entrada}
\KwOut{Datos de salida}

\For{cada elemento en los datos de entrada}{
    Procesar elemento\;
    \If{condición}{
        Ejecutar acción A\;
    }
    \Else{
        Ejecutar acción B\;
    }
}
\Return{Resultado}\;
\end{algorithm}

Los algoritmos deben estar numerados y seguir las mismas reglas que las figuras: deben ser referenciados en el texto antes de su aparición y su leyenda debe contener punto final.